\definecolor{headergreen}{RGB}{90, 164, 90}
\definecolor{tablerowgreen}{RGB}{235, 247, 238}
\definecolor{tablerowwhite}{RGB}{255, 255, 255}

% \omiss produces '[...]'
\newcommand{\omissis}{[\dots\negthinspace]}

% rimuove la scritta capitolo X
\titleformat{\chapter}[display]{\normalfont\huge\bfseries}{}{0pt}{\huge}
\titlespacing*{\chapter}{0pt}{-20pt}{20pt}

% Itemize symbols
% see: https://tex.stackexchange.com/a/62497
% \renewcommand{\labelitemi}{$\bullet$}
% \renewcommand{\labelitemii}{$\cdot$}
% \renewcommand{\labelitemiii}{$\diamond$}
% \renewcommand{\labelitemiv}{$\ast$}


\let\Chaptermark\chaptermark
% \Chaptername gives current chapter name
\def\chaptermark#1{\def\Chaptername{#1}\Chaptermark{#1}}
\makeatletter
% \currentname gives the current section name
\newcommand*{\currentname}{\@currentlabelname}
\makeatother

% Uncomment the following line for a different header/footer style
\pagestyle{fancy} \setlength{\headheight}{14.5pt}
\fancyhead[L]{\fontsize{12}{14.5} \selectfont \thechapter. \Chaptername}
\fancyhead[R]{\fontsize{12}{14.5} \selectfont \currentname}

\renewcommand{\arraystretch}{1.8}
\newcolumntype{L}[1]{>{\raggedright\arraybackslash}p{#1}}
\newcommand{\thead}[1]{\textcolor{white}{\textls[110]{\textbf{#1}}}}

% Page number always in footer
\cfoot{\thepage}


% Custom hyphenation rules
\hyphenation {
    e-sem-pio
    ex-am-ple
}

% Images path, not using \graphicspath because it doesn't properly work with
% latexmk custom dependencies
\NewCommandCopy{\latexincludegraphics}{\includegraphics}
\renewcommand{\includegraphics}[2][]{\latexincludegraphics[#1]{../images/#2}}

% Page format settings
% see: http://wwwcdf.pd.infn.it/AppuntiLinux/a2547.htm
\setlength{\parindent}{0pt}    % first row indentation
\setlength{\parskip}{0pt}       % paragraphs spacing


% Load variables
\newcommand{\myName}{Angela Canazza}
\newcommand{\myID}{2111030}
\newcommand{\myTitle}{AI Document Consistency Platform – Controllo qualità
documentale automatizzato}
\newcommand{\myDegree}{Tesi di laurea triennale}
\newcommand{\myUni}{Università degli Studi di Padova}
% For BSc level just use "Corso di Laurea" and don't add "Triennale" to it
\newcommand{\myFaculty}{Corso di Laurea in Informatica}
\newcommand{\myDepartment}{Dipartimento di Matematica ``Tullio Levi-Civita''}
\newcommand{\profTitle}{Prof.}
\newcommand{\myProf}{Marco Zanella}
\newcommand{\myLocation}{Padova}
\newcommand{\myAA}{2025-2026}
\newcommand{\myTime}{Luglio 2026}

% PDF file metadata fields
% when updating them delete the build directory, otherwise they won't change
\begin{filecontents*}{\jobname.xmpdata}
  \Title{Document's title}
  \Author{Author's name}
  \Language{it-IT}
  \Subject{Short description}
  \Keywords{keyword1\sep keyword2\sep keyword3}
\end{filecontents*}


% Acronyms
\newacronym[description={\glslink{api}{Application Program Interface}}]
    {apig}{APIg}{Application Programming Interface}

\newacronym[description={\glslink{llm}{Large Language Model}}]
    {llmg}{LLMg}{Large Language Model}

\newacronym[description={\glslink{ai}{Artificial Intelligence}}]
    {aig}{AIg}{Artificial Intelligence}

\newacronym[description={\glslink{jwt}{JSON Web Token}}]
    {jwtg}{JWTg}{JSON Web Token}

% \newacronym[description={\glslink{umlg}{Unified Modeling Language}}]
%     {uml}{UML}{Unified Modeling Language}

% Glossary entries
\newglossaryentry{api} {
    name=\glslink{api}{API},
    text=API,
    sort=api,
    description={in informatica con il termine \emph{API} (ing. \emph{Application Programming Interface}, interfaccia di programmazione di un'applicazione) si indica un insieme di definizioni e protocolli che permettono a due applicazioni software di comunicare tra loro. Le \emph{API} definiscono le modalità con cui un componente software può richiedere servizi o dati a un altro, astraendo i dettagli implementativi interni e rendendo possibile l'integrazione tra sistemi eterogenei}
}

% \newglossaryentry{agile} {
%     name=\glslink{agile}{Agile},
%     text=Agile,
%     sort=agile,
%     description={in ingegneria del software con il termine \emph{Agile} si indica un insieme di metodologie di sviluppo software basate su cicli iterativi e incrementali. L'approccio privilegia la collaborazione continua tra i membri del team, con il cliente e la capacità di adattarsi ai cambiamenti dei requisiti}
% }

\newglossaryentry{git} {
    name=\glslink{git}{Git},
    text=Git,
    sort=git,
    description={Git è un sistema di controllo di versione distribuito, creato nel 2005. Permette di tracciare le modifiche al codice sorgente nel tempo, gestire branch paralleli di sviluppo e unire le modifiche tramite operazioni di merge. A differenza dei sistemi centralizzati, ogni sviluppatore dispone di una copia completa della storia del repository in locale}
}

\newglossaryentry{github} {
    name=\glslink{github}{GitHub},
    text=GitHub,
    sort=github,
    description={GitHub è una piattaforma web per l'hosting e la gestione di repository Git. Permette il versionamento del codice sorgente, la collaborazione tra sviluppatori tramite pull request e la gestione delle modifiche attraverso un sistema di branch. È uno degli strumenti di versionamento più diffusi nel settore dello sviluppo software}
}

\newglossaryentry{llm} {
    name=\glslink{llm}{LLM},
    text=Large Language Model,
    sort=llm,
    description={in informatica con il termine \emph{Large Language Model} (LLM, ing. modello linguistico di grandi dimensioni) si indica un tipo avanzato di modello di intelligenza artificiale generativa, addestrato su enormi quantità di testo con lo scopo di comprendere, elaborare e generare linguaggio naturale o codice. Questi modelli sono in grado di svolgere compiti complessi quali sintesi, traduzione, risposta a domande e analisi semantica del testo, senza necessità di essere ri-addestrati per ciascun compito specifico}
}

\newglossaryentry{jira} {
    name=\glslink{jira}{Jira},
    text=Jira,
    sort=jira,
    description={Jira è uno strumento software sviluppato da Atlassian per la gestione di progetti e il tracciamento delle attività. Consente di organizzare il lavoro in task, assegnarle ai membri del team, definirne le priorità e monitorarne l'avanzamento. È ampiamente adottato in ambito informatico per supportare metodologie di sviluppo agili, in particolare attraverso la gestione di sprint e backlog}
}

\newglossaryentry{docker} {
    name=\glslink{docker}{Docker},
    text=Docker,
    sort=docker,
    description={Docker è una piattaforma \emph{open source} per la creazione, la distribuzione e l'esecuzione di applicazioni all'interno di \emph{container}.
    Un \emph{container} è un'unità software che racchiude il codice dell'applicazione
    insieme a tutte le sue dipendenze, garantendo che il programma funzioni in modo uniforme indipendentemente dall'ambiente di esecuzione}
}

\newglossaryentry{swagger} {
    name=\glslink{swagger}{Swagger},
    text=Swagger,
    sort=swagger,
    description={Swagger è un insieme di strumenti \emph{open source} basati sulla specifica \emph{OpenAPI}, utilizzati per progettare, documentare e testare \emph{API} REST. Permette di generare automaticamente una interfaccia web interattiva che descrive gli \emph{endpoint} disponibili, i parametri accettati e i formati di risposta, facilitando l'integrazione tra sistemi diversi e la comunicazione tra sviluppatori}
}

\newglossaryentry{slack} {
    name=\glslink{slack}{Slack},
    text=Slack,
    sort=slack,
    description={Slack è una piattaforma di comunicazione aziendale che consente di organizzare le conversazioni in canali tematici, inviare messaggi diretti e condividere file tra i membri di un team. È ampiamente adottata in ambito informatico come strumento di comunicazione interna nei progetti di sviluppo software}
}

\newglossaryentry{ai} {
    name=\glslink{ai}{AI},
    text=AI,
    sort=ai,
    description={Con il termine \emph{AI} (ing. \emph{Artificial Intelligence}, intelligenza artificiale) si indica un ramo dell'informatica che studia e sviluppa sistemi in grado di svolgere compiti che richiederebbero normalmente l'intelligenza umana, quali il riconoscimento del linguaggio naturale, l'apprendimento automatico, il ragionamento e la risoluzione di problemi. Negli ultimi anni il termine è spesso associato in particolare ai modelli di \emph{machine learning} e ai \gls{llm}}
}

\newglossaryentry{rulebased} {
    name=\glslink{rulebased}{Rule-Based},
    text=Rule-Based,
    sort=rulebased,
    description={Con il termine \emph{Rule-Based} si indica un approccio algoritmico in cui il comportamento del sistema è definito da un insieme di regole esplicite e deterministiche. Un sistema \emph{rule-based} produce sempre lo stesso output a parità di input, garantendo risultati prevedibili e facilmente testabili}
}

\newglossaryentry{rendanheyi} {
    name=\glslink{rendanheyi}{RenDanHeYi},
    text=RenDanHeYi,
    sort=rendanheyi,
    description={il \emph{RenDanHeYi} è un modello organizzativo sviluppato da Zhang Ruimin, fondatore del gruppo Haier, nel 2005. Il termine è la contrazione delle parole cinesi per "persona", "ordine" e "unificazione". Rappresenta una visione in cui il rapporto tra impresa e individuo non è gerarchico ma reciproco: l'azienda mette le persone in condizione di esprimere il proprio potenziale, che a sua volta genera valore per l'organizzazione. Le unità operative, dette micro-imprese, godono di autonomia decisionale e responsabilità diretta sui risultati.}
}

\newglossaryentry{holacracy} {
    name=\glslink{holacracy}{Holacracy},
    text=Holacracy,
    sort=holacracy,
    description={\emph{Holacracy} è una tecnologia sociale per la gestione organizzativa che redistribuisce il potere decisionale all'interno dell'azienda secondo ruoli dinamici e processi strutturati, in contrapposizione alla tradizionale gerarchia manageriale. Il termine è stato coniato da Brian Robertson nel 2007..}
}

\newglossaryentry{openplatform} {
    name=\glslink{openplatform}{Open Platform Organization},
    text=Open Platform Organization,
    sort=openplatformorganization,
    description={\emph{Open Platform Organization} è un modello organizzativo che integra i principi del rendanheyi e dell'holacracy per creare un'organizzazione in cui la \emph{leadership} è diffusa, le decisioni sono distribuite e le diverse funzioni aziendali collaborano attorno a uno scopo comune.}
}

\newglossaryentry{multer} {
    name=\glslink{multer}{Multer},
    text=Multer,
    sort=multer,
    description={\emph{Multer} è un middleware per Node.js progettato per la gestione delle richieste HTTP di tipo \emph{multipart/form-data}, utilizzato principalmente per la ricezione di file caricati dal client. Espone il contenuto del file come buffer in memoria, rendendolo disponibile al livello applicativo senza richiedere la scrittura su disco.}
}

\newglossaryentry{mammoth} {
    name=\glslink{mammoth}{Mammoth},
    text=Mammoth,
    sort=mammoth,
    description={\emph{Mammoth} è una libreria per Node.js che converte documenti in formato DOCX in HTML, preservando la struttura semantica del documento originale come titoli, paragrafi, elenchi e tabelle.}
}

\newglossaryentry{turndown} {
    name=\glslink{turndown}{Turndown},
    text=Turndown,
    sort=turndown,
    description={\emph{Turndown} è una libreria per Node.js che converte HTML in Markdown. Supporta l'estensione tramite plugin, tra cui \emph{turndown-plugin-gfm} che aggiunge il supporto alla sintassi GitHub Flavored Markdown, necessaria per la corretta rappresentazione delle tabelle.}
}

\newglossaryentry{docling} {
    name=\glslink{docling}{Docling},
    text=Docling,
    sort=docling,
    description={\emph{Docling} è un progetto open source sviluppato da IBM Research per l'estrazione e la conversione di documenti PDF in formati strutturati come Markdown. A differenza degli estrattori di testo tradizionali, Docling tratta ogni pagina come un'immagine e applica modelli di machine learning per riconoscere la struttura visiva del documento prima di estrarne il contenuto, consentendo di preservare titoli, tabelle ed elenchi nel risultato finale. L'integrazione con backend non Python avviene tramite Docling Serve, un server REST che espone le funzionalità della libreria via HTTP.}
}

\newglossaryentry{strategy} {
    name=\glslink{strategy}{Strategy},
    text=Strategy,
    sort=strategy,
    description={\emph{Strategy} è un pattern di progettazione comportamentale che definisce una famiglia di algoritmi, li incapsula ciascuno in una classe separata e li rende intercambiabili attraverso un'interfaccia comune. Il chiamante non conosce la strategia concreta utilizzata, ma interagisce unicamente con l'interfaccia, delegando l'implementazione specifica alla strategia selezionata. Questo approccio favorisce l'estensibilità: aggiungere un nuovo algoritmo richiede unicamente la creazione di una nuova classe aderente all'interfaccia, senza modificare il codice esistente.}
}

\newglossaryentry{cognito} {
    name=\glslink{cognito}{AWS Cognito},
    text=AWS Cognito,
    sort=cognito,
    description={\emph{AWS Cognito} è un servizio di identity management fornito da Amazon Web Services che gestisce la registrazione, l'autenticazione e il controllo degli accessi degli utenti. Permette di definire gruppi di utenti a cui assegnare ruoli distinti e si integra con le applicazioni tramite standard aperti come OAuth 2.0 e OpenID Connect, emettendo token JWT alla conclusione del flusso di autenticazione.}
}

\newglossaryentry{jwt} {
    name=\glslink{jwt}{JWT},
    text=JWT,
    sort=jwt,
    description={\emph{JWT} (JSON Web Token) è uno standard aperto che definisce un formato compatto per la trasmissione sicura di informazioni tra due parti sotto forma di oggetto JSON. Un JWT è composto da tre parti codificate in Base64 e separate da un punto: l'header, che specifica il tipo di token e l'algoritmo di firma; il payload, che contiene le informazioni sull'utente dette \emph{claims}; e la firma, che garantisce l'integrità del token e ne permette la verifica tramite chiave pubblica senza dover contattare il server emittente.}
}

\newglossaryentry{amplify} {
    name=\glslink{amplify}{AWS Amplify},
    text=AWS Amplify,
    sort=amplify,
    description={\emph{AWS Amplify} è una libreria sviluppata da Amazon Web Services che astrae le chiamate dirette ai servizi AWS lato frontend. Nel contesto dell'autenticazione, si integra con AWS Cognito gestendo il flusso di login, la ricezione e il rinnovo dei token e l'accesso alla sessione corrente dell'utente tramite metodi come \texttt{fetchAuthSession}.}
}

\newglossaryentry{interceptor} {
    name=\glslink{interceptor}{Interceptor},
    text=Interceptor,
    sort=interceptor,
    description={\emph{Interceptor} è un meccanismo che permette di intercettare 
    le richieste HTTP in uscita o le risposte in entrata prima che vengano 
    elaborate dall'applicazione, consentendo di modificarle in modo centralizzato. 
    Nel contesto di Axios, viene utilizzato per aggiungere automaticamente 
    l'header di autenticazione a ogni richiesta senza dover ripetere la logica 
    in ogni chiamata.}
}

% \newglossaryentry{umlg} {
%     name=\glslink{uml}{UML},
%     text=UML,
%     sort=uml,
%     description={in ingegneria del software \emph{UML, Unified Modeling Language} (ing. linguaggio di modellazione unificato) è un linguaggio di modellazione e specifica basato sul paradigma object-oriented. L'\emph{UML} svolge un'importantissima funzione di ``lingua franca'' nella comunità della progettazione e programmazione a oggetti. Gran parte della letteratura di settore usa tale linguaggio per descrivere soluzioni analitiche e progettuali in modo sintetico e comprensibile a un vasto pubblico}
% }

\makeglossaries

\bibliography{appendix/bibliography}

\defbibheading{bibliography} {
    \cleardoublepage
    \phantomsection
    \addcontentsline{toc}{chapter}{\bibname}
    \chapter*{\bibname\markboth{\bibname}{\bibname}}
}

% Spacing between entries
\setlength\bibitemsep{1.5\itemsep}

\DeclareBibliographyCategory{opere}
\DeclareBibliographyCategory{web}

\addtocategory{opere}{womak:lean-thinking}
\addtocategory{web}{site:agile-manifesto}

\defbibheading{opere}{\section*{Riferimenti bibliografici}}
\defbibheading{web}{\section*{Siti Web consultati}}


\captionsetup{
    tableposition=top,
    figureposition=bottom,
    font=small,
    format=hang,
    labelfont=bf
}

\hypersetup{
    %hyperfootnotes=false,
    %pdfpagelabels,
    colorlinks=true,
    linktocpage=true,
    pdfstartpage=1,
    pdfstartview=,
    breaklinks=true,
    pdfpagemode=UseNone,
    pageanchor=true,
    pdfpagemode=UseOutlines,
    plainpages=false,
    bookmarksnumbered,
    bookmarksopen=true,
    bookmarksopenlevel=1,
    hypertexnames=true,
    pdfhighlight=/O,
    %nesting=true,
    %frenchlinks,
    urlcolor=webbrown,
    linkcolor=OliveGreen,
    citecolor=OliveGreen  
    %pagecolor=RoyalBlue,
}

% Delete all links and show them in black
\if \isprintable 1
    \hypersetup{draft}
\fi

% Listings setup
\lstset{
    language=[LaTeX]Tex,%C++,
    keywordstyle=\color{RoyalBlue}, %\bfseries,
    basicstyle=\small\ttfamily,
    %identifierstyle=\color{NavyBlue},
    commentstyle=\color{Green}\ttfamily,
    stringstyle=\rmfamily,
    numbers=none, %left,%
    numberstyle=\scriptsize, %\tiny
    stepnumber=5,
    numbersep=8pt,
    showstringspaces=false,
    breaklines=true,
    frameround=ftff,
    frame=single
}

\definecolor{webgreen}{rgb}{0,.5,0}
\definecolor{webbrown}{rgb}{.6,0,0}

\newcommand{\sectionname}{sezione}
\addto\captionsitalian{\renewcommand{\figurename}{Figura}
                       \renewcommand{\tablename}{Tabella}}

\newcommand{\glsfirstoccur}{\ap{{[g]}}}

\newcommand{\intro}[1]{\emph{\textsf{#1}}}

% Risks environment
\newcounter{riskcounter}                % define a counter
\setcounter{riskcounter}{0}             % set the counter to some initial value

%%%% Parameters
% #1: Title
\newenvironment{risk}[1]{
    \refstepcounter{riskcounter}        % increment counter
    \par \noindent                      % start new paragraph
    \textbf{\arabic{riskcounter}. #1}   % display the title before the content of the environment is displayed
}{
    \par\medskip
}

\newcommand{\riskname}{Rischio}

\newcommand{\riskdescription}[1]{\textbf{\\Descrizione:} #1.}

\newcommand{\risksolution}[1]{\textbf{\\Soluzione:} #1.}

% Use case environment
\newcounter{usecasecounter}             % define a counter
\setcounter{usecasecounter}{0}          % set the counter to some initial value

%%%% Parameters
% #1: ID
% #2: Nome
\newenvironment{usecase}[2]{
    \renewcommand{\theusecasecounter}{\usecasename #1}  % this is where the display of
                                                        % the counter is overwritten/modified
    \refstepcounter{usecasecounter}             % increment counter
    \vspace{10pt}
    \par \noindent                              % start new paragraph
    {\large \textbf{\usecasename #1: #2}}       % display the title before the
                                                % content of the environment is displayed
    \medskip
}{
    \medskip
}

\newcommand{\usecasename}{UC}

\newcommand{\usecaseactors}[1]{\textbf{\\Attori principali:} #1. \vspace{4pt}}
\newcommand{\usecasesecondactors}[1]{\textbf{\\Attori secondari:} #1. \vspace{4pt}}
\newcommand{\usecasepre}[1]{\textbf{\\Precondizioni:} #1. \vspace{4pt}}
\newcommand{\usecasedesc}[1]{\textbf{\\Descrizione:} #1. \vspace{4pt}}
\newcommand{\usecasepost}[1]{\textbf{\\Postcondizioni:} #1 \vspace{4pt}}
\newcommand{\usecasealt}[1]{\textbf{\\Scenario alternativo:} #1 \vspace{4pt}}
\newcommand{\usecaseprin}[1]{\textbf{\\Scenario principale:} #1 \vspace{4pt}}
\newcommand{\usecasetrig}[1]{\textbf{\\Trigger:} #1. \vspace{4pt}}
\newcommand{\usecaseincl}[1]{\textbf{\\Inclusioni:} #1. \vspace{4pt}}
\newcommand{\usecaseext}[1]{\textbf{\\Estensioni:} #1. \vspace{4pt}}
\newcommand{\usecasegen}[1]{\textbf{\\Generalizzazioni:} #1. \vspace{4pt}}

% Namespace description environment
\newenvironment{namespacedesc}{
    \vspace{10pt}
    \par \noindent  % start new paragraph
    \begin{description}
}{
    \end{description}
    \medskip
}

\newcommand{\classdesc}[2]{\item[\textbf{#1:}] #2}
\newcommand{\ucref}[1]{\hyperref[uc:#1]{[UC#1]}}
\definecolor{tablegray}{RGB}{180, 200, 185}
\setlength{\arrayrulewidth}{0.2pt}
\arrayrulecolor{tablegray}